\section{Problem Formulation}
\label{sec:problem_formulation}

We consider a multi-agent reasoning system built on large language models (LLMs).
Given an input query \(x\), the objective is to produce a correct final answer \(y\) through a bounded number of interaction rounds among an author agent, \(k\) reviewer agents, and a meta-reviewer agent.

Existing multi-agent LLM frameworks typically represent the reasoning process implicitly, as a sequence of natural-language messages exchanged among agents~\cite{du2024improving,qian2024chatdev,hong2024metagpt,wang2025mars}.
Such message-based representations do not expose an explicit object that records which candidate solution is currently under consideration, how it was derived, and how reliable it is judged to be.
We therefore formulate multi-agent reasoning as a \emph{state-transition process} over an abstract task environment, inspired by world-model approaches to sequential decision making~\cite{ha2018recurrent,hafner2020dream}.

\paragraph{States, actions, transitions.}
Let \(s_t \in \mathcal{S}\) denote the task state after the \(t\)-th committed action, with \(s_0\) derived from the query \(x\).
An action \(a_t \in \mathcal{A}\) is a typed proposal produced by an agent (a draft or a revision of a candidate solution).
Committing an action yields a new state through a transition function
\begin{equation}
s_{t+1} = T(s_t, a_t, o_t),
\end{equation}
where \(o_t\) is an optional observation (in our setting, the index of the review round).
A reasoning \emph{trajectory} is the sequence \(\xi = (s_0, s_1, \ldots, s_H)\) with horizon \(H\).

\paragraph{Validation and risk.}
Two evaluation functions operate on states. A validator
\begin{equation}
V(s_t, a_t) \rightarrow (\mathit{ok}, \mathit{hint})
\end{equation}
decides whether a proposed action may be passed on to the reviewers, and a rollout function
\begin{equation}
R(s_t, \mathcal{F}_t) \rightarrow r_t \in [0,1]
\end{equation}
maps the current state and the reviewer feedback \(\mathcal{F}_t\) collected for it to a \emph{trajectory risk}: an estimate of how likely it is that accepting the current candidate would be a mistake.

\paragraph{Coordination as a policy.}
A coordination mechanism \(\pi\) specifies who owns the state, which information each agent sees, and under which condition a candidate is accepted.
We study two mechanisms (centralized and peer-to-peer, Section~\ref{sec:method}) and an uncertainty-dependent rule for choosing between them: for a task distribution \(D\) with coordination uncertainty \(U(D)\),
\begin{equation}
\label{eq:policy}
\pi^{*}(D) =
\begin{cases}
\text{Centralized}, & U(D) < \tau,\\
\text{Peer}, & U(D) \ge \tau,
\end{cases}
\end{equation}
where \(\tau\) is a threshold.
Eq.~\eqref{eq:policy} is a hypothesis to be tested empirically, not an optimality result; Section~\ref{sec:cui} defines \(U(D)\) and Section~\ref{sec:cui_results} examines how well the rule holds.

\section{World-Model--Guided Coordination}
\label{sec:method}

Our framework consists of (i) a world model that stores and evaluates reasoning states, (ii) LLM agents in three roles, and (iii) a coordination mechanism that determines how the agents use the world model.
This section specifies each component at the level of detail needed to re-implement it.
The world model is deliberately lightweight: it is a symbolic state store with rule-based evaluation functions and involves no additional LLM calls and no learned parameters.

\subsection{World-Model State and Operations}
\label{sec:wm}

\subsubsection{State composition}
The world model maintains a rooted tree of state nodes. Each node is a tuple
\begin{equation}
s = \langle \mathit{id},\ \mathit{parent},\ a,\ o \rangle ,
\end{equation}
where \(\mathit{id}\) is a unique pointer, \(\mathit{parent}\) is the pointer of the preceding state (empty for \(s_0\)), \(a = \langle \mathit{type}, \mathit{content}, \mathit{meta} \rangle\) is the action that produced the state, and \(o\) is the observation attached to it.
The action type is one of \textsc{init} (the query), \textsc{draft} (the author's first candidate), \textsc{revision} (a candidate revised after meta-review feedback, whose \(\mathit{meta}\) field stores that feedback), and \textsc{revalidated} (a candidate re-submitted after a failed validation).
The \(\mathit{content}\) field holds the full text of the candidate solution, including its reasoning steps and final answer.
The path from the root to the current node is therefore a complete, ordered record of every candidate that has been considered for the query and of the feedback that triggered each revision.
Agents never exchange this structure itself; they hold the pointer of the current node and query the world model through the four operations below.

\subsubsection{Transition \(T\)}
\(T(s_t, a_t, o_t)\) allocates a new node whose parent is \(s_t\), stores \(a_t\) and \(o_t\) in it, and returns its pointer. States are immutable once created, so a revision never overwrites the candidate it replaces.

\subsubsection{Validation \(V\)}
\(V(s_t, a_t)\) applies two rule-based checks to the content of a proposed action:
(i) a completeness check, which rejects contents shorter than 10 characters; and
(ii) a degenerate-result check, which rejects contents that match a small lexicon \(\mathcal{E}\) of failure indicators (a final expression that evaluates to zero, or the tokens ``undefined'', ``infinity'', and ``cannot'').
If a check fails, the validator returns \(\mathit{ok}=0\) together with a textual hint; the author is then asked once to re-check and re-submit the candidate in the required format, and the re-submitted candidate is committed as a \textsc{revalidated} state.
The implementation also contains an optional table of per-pattern outcome statistics that can be updated from past results. We disable it in all experiments: every query starts from an empty world model, so no ground-truth information from other test questions can influence a run.

\subsubsection{State summary \(\sigma\)}
\(\sigma(s_t)\) renders the current state for consumption by other agents. It returns the content of the latest action truncated to 800 characters.
In centralized mode, a state whose content matches \(\mathcal{E}\) is not forwarded verbatim; \(\sigma\) instead returns a short warning that the previous step has potential issues.
The summary is appended to each reviewer prompt in a dedicated \texttt{[WORLD\_STATE]} block.

\subsubsection{Rollout risk \(R\)}
Let \(\mathcal{F}_t=\{f_1,\ldots,f_k\}\) be the reviews of the current candidate in review round \(n\in\{1,2,\ldots\}\).
Each review \(f_i\) carries a reliability score \(c_i\in[0,1]\), computed as the exponential of the mean token log-probability of the reviewer's response, and a polarity \(u_i\in\{+1,0,-1\}\), obtained by comparing the number of matches against a positive lexicon (e.g., ``correct'', ``accurate'', ``agree'') and a negative lexicon (e.g., ``incorrect'', ``error'', ``mistake'').
Let \(\rho = \tfrac{1}{k}\sum_i \mathbb{1}[u_i=+1]\) be the fraction of positive reviews and \(\bar c = \tfrac1k \sum_i c_i\) the mean reliability.
The trajectory risk is
\begin{equation}
\label{eq:risk}
r_t = \mathrm{clip}_{[0,1]}\Big( b(n,\delta_t) + 0.3\,e_t + 0.25\,(1-\rho) + g(\rho) + h(\bar c) \Big),
\end{equation}
where \(\delta_t\) is the depth of \(s_t\) in the state tree, \(e_t=\mathbb{1}[\mathit{content}(s_t)\text{ matches }\mathcal{E}]\), and
\begin{align*}
b(n,\delta) &= \begin{cases} 0.30, & n = 1,\\ 0.20 + 0.01\,\delta, & n \ge 2, \end{cases}\\
g(\rho) &= \begin{cases} +0.10, & \rho < 0.5,\\ +0.05, & 0.5 \le \rho < 0.7,\\ +0.02, & 0.7 \le \rho < 1,\\ -0.10, & \rho = 1, \end{cases}\\
h(\bar c) &= \begin{cases} +0.20, & \bar c < 0.6,\\ +0.10, & 0.6 \le \bar c < 0.7,\\ -0.05, & 0.7 \le \bar c \le 0.85,\\ -0.15, & \bar c > 0.85. \end{cases}
\end{align*}
The risk thus rises with reviewer disagreement, low reviewer reliability, and degenerate content, and its base level depends on how far the trajectory has progressed.
The constants in Eq.~\eqref{eq:risk} were set manually during development and are kept fixed across all datasets and backbones.

\subsection{Agent Roles}

\paragraph{Author} The author generates a candidate solution in a fixed ``Thoughts / Answer'' format and revises it when feedback is returned.

\paragraph{Reviewers} Each of the \(k\) reviewers first solves the problem independently, then compares its result with the author's candidate and returns a decision (right/wrong), a justification, and its own final answer. Reviewers do not see each other's reviews. When a world model is used, the reviewer prompt additionally contains the state summary \(\sigma(s_t)\).

\paragraph{Meta-reviewer} The meta-reviewer receives the candidate, all reviews with their reliability scores, and (when a world model is used) the risk estimate \(r_t\). It returns a decision (right/wrong), a justification, and suggestions for revision.

\subsection{Coordination Architectures}

\subsubsection{Centralized coordination}
In the centralized architecture (Fig.~\ref{fig:centralized}) a single world-model instance is the only holder of the reasoning state.
The author submits each candidate to it as an action; the world model validates the action, commits the transition, and hands the summary \(\sigma(s_t)\) to every reviewer, so all reviewers evaluate exactly the same committed state.
After the reviews are collected, the world model computes \(r_t\) with Eq.~\eqref{eq:risk}, passes it to the meta-reviewer, and finally applies the acceptance gate of Section~\ref{sec:gate}.
A candidate can therefore only be returned as the final answer if the world model agrees; a positive meta-review alone is not sufficient.

\begin{figure*}[!t]
    \centering
    \resizebox{0.95\textwidth}{!}{\begin{tikzpicture}[node distance=7mm and 12mm]
\node[agent] (q) {Query $x$};
\node[agent, right=8mm of q, minimum width=16mm] (a) {Author};
\node[wm, right=22mm of a, minimum width=68mm, minimum height=30mm] (w) {};
\node[ttl, anchor=north] at ([yshift=-1mm]w.north) {Shared world model (single state store)};
\node[st] (s0) at ([xshift=-20mm,yshift=2mm]w.center) {$s_0$};
\node[st] (s1) at ([xshift=-7mm,yshift=2mm]w.center) {$s_1$};
\node[font=\scriptsize] (sd) at ([xshift=6mm,yshift=2mm]w.center) {$\cdots$};
\node[st] (s2) at ([xshift=19mm,yshift=2mm]w.center) {$s_t$};
\draw[flow] (s0) -- (s1); \draw[flow] (s1) -- (sd); \draw[flow] (sd) -- (s2);
\node[op] (v) at ([xshift=-22mm,yshift=-9.5mm]w.center) {$V$: validate};
\node[op] (sg) at ([xshift=0mm,yshift=-9.5mm]w.center) {$\sigma$: summarize};
\node[op] (rr) at ([xshift=22.5mm,yshift=-9.5mm]w.center) {$R$: rollout risk};
\node[agent, right=24mm of w, yshift=8mm, minimum width=20mm] (r1) {Reviewer 1};
\node[agent, below=3mm of r1, minimum width=20mm] (r2) {Reviewer $k$};
\node[gate, below=13mm of w] (g) {WM-guided\\acceptance gate};
\node[agent, minimum width=20mm] (m) at (r2 |- g) {Meta-reviewer};
\node[agent, below=8mm of g, fill=green!10] (out) {Final answer $y$};

\draw[flow] (q) -- (a);
\draw[flow] ([yshift=2mm]a.east) -- node[lbl, above] {\ding{172} action $a_t$} ([yshift=2mm]a.east -| w.west);
\draw[fb] ([yshift=-2mm]a.east -| w.west) -- node[lbl, below] {\ding{173} invalid: revise} ([yshift=-2mm]a.east);
\draw[flow] ([yshift=8mm]w.east) -- node[lbl, above] {\ding{174} $\sigma(s_t)$} (r1.west);
\draw[flow] ([yshift=8mm]w.east) -- ([yshift=1.5mm]r2.west);
\draw[flow] ([yshift=-2mm]r2.west) -- node[lbl, below, pos=0.45] {\ding{175} reviews} ([yshift=-2mm]r2.west -| w.east);
\draw[flow] (r2) -- node[lbl, right] {\ding{175} reviews} (m);
\draw[flow] ([xshift=27mm]w.south) |- node[lbl, pos=0.25, right] {\ding{176} risk $r_t$} ([yshift=2mm]m.west);
\draw[flow] ([yshift=-2mm]m.west) -- node[lbl, below, pos=0.5] {\ding{177} decision} ([yshift=-2mm]m.west -| g.east) -- (g.east);
\draw[flow] (w.south) -- node[lbl, left] {$r_t$, consensus,\\confidence} (g.north);
\draw[flow, green!45!black] (g) -- node[lbl, right] {accept} (out);
\draw[fb] (g.west) -| node[lbl, pos=0.25, above] {\ding{178} reject: feedback to author;}  node[lbl, pos=0.25, below] {revision $a_{t+1}$, $s_{t+1}{=}T(s_t,a_{t+1})$} (a.south);
\end{tikzpicture}
}
    \caption{Centralized world-model coordination. All agents read from and write to one shared state store; the world model validates each action, distributes the state summary, computes the trajectory risk from the reviews, and gates acceptance.}
    \label{fig:centralized}
\end{figure*}

\subsubsection{Peer-to-peer coordination}
In the peer architecture (Fig.~\ref{fig:peer}) there is no shared state store. Every agent owns a local world-model instance.
The author commits and validates its candidates in its local instance and sends the resulting summary to the reviewers as part of its message; each reviewer and the meta-reviewer mirror the received information into their own local instances.
The meta-reviewer's local world model has no access to a shared review record and therefore computes only a local risk (Eq.~\eqref{eq:risk} without the review-dependent terms, with base level 0.5), which is reported in the meta-reviewer prompt.
There is no global gate: a candidate is accepted as soon as the meta-reviewer decides ``right''.
Peer coordination is thus the lighter mechanism. It keeps explicit state tracking and validation on the author side but leaves the acceptance decision to the agents, whereas centralized coordination adds a conservative, state-level check that can overrule them.

\begin{figure*}[!t]
    \centering
    \resizebox{0.95\textwidth}{!}{\begin{tikzpicture}[node distance=7mm and 12mm]
\node[agent] (q) {Query $x$};
\node[agent, right=8mm of q, minimum width=22mm] (a) {Author};
\node[wm, below=1.2mm of a, minimum width=22mm, font=\scriptsize] (wa) {local WM$_A$: $T$, $V$, $\sigma$};
\node[agent, right=62mm of a, yshift=8mm, minimum width=20mm] (r1) {Reviewer 1};
\node[wm, right=1.2mm of r1, font=\scriptsize] (w1) {local WM$_1$};
\node[agent, below=3mm of r1, minimum width=20mm] (r2) {Reviewer $k$};
\node[wm, right=1.2mm of r2, font=\scriptsize] (w2) {local WM$_k$};
\node[agent, below=15mm of r2, minimum width=20mm] (m) {Meta-reviewer};
\node[wm, right=1.2mm of m, font=\scriptsize] (wmm) {local WM$_M$: $R$};
\node[agent, fill=green!10] (out) at ($(a |- m)!0.3!(m)$) {Final answer $y$};
\node[font=\scriptsize, align=center, text width=44mm] at ($(a.east)!0.5!(r2.west)+(0,-11mm)$) {no shared state store,\\no global acceptance gate};

\draw[flow] (q) -- (a);
\draw[flow] (a.east) -- node[lbl, above, sloped] {\ding{172} answer + summary $\sigma(s^{A}_t)$} (r1.west);
\draw[flow] (a.east) -- (r2.west);
\draw[flow] (r2) -- node[lbl, right] {\ding{173} reviews} (m);
\draw[flow, green!45!black] (m.west) -- node[lbl, above] {\ding{174} ``right'': accept} (out.east);
\draw[fb] (m.south) -- ++(0,-5mm) -| node[lbl, pos=0.25, below] {\ding{175} ``wrong'': feedback, author revises and re-validates locally} (wa.south);
\end{tikzpicture}
}
    \caption{Peer-to-peer world-model coordination. Each agent keeps a local world-model replica; state summaries travel inside messages, and acceptance follows the meta-reviewer's decision.}
    \label{fig:peer}
\end{figure*}

\subsection{World-Model--Guided Acceptance}
\label{sec:gate}

In centralized mode a candidate that the meta-reviewer judges ``right'' is accepted only if it also passes the following gate.
Let \(\kappa\) be the number of reviewers counted as agreeing with the candidate (a reviewer agrees if its review has positive polarity or its reliability exceeds 0.7).
The risk is first adjusted by lexical cues in the meta-review, \(r'_t = \mathrm{clip}_{[0,1]}(r_t + \Delta_{\mathrm{meta}})\), where \(\Delta_{\mathrm{meta}}\in[-0.15,+0.25]\) is positive when the meta-review mentions flaws or concerns and negative when it states that the answer is clearly correct.
The \emph{effective risk} combines the adjusted risk with reviewer reliability,
\begin{equation}
\tilde r_t = 0.7\, r'_t + 0.3\,(1-\bar c),
\end{equation}
and is compared with a length-aware threshold \(\theta = \theta_0 + \min(0.12,\ L/900)\), where \(L\) is the number of words in the candidate.
The candidate is accepted if either
\begin{align}
&\tilde r_t \le \theta \ \wedge\ \kappa \ge \kappa_{\min} \ \wedge\ \bar c > c_{\min}, \quad\text{or} \label{eq:gateA}\\
&\kappa = k \ \wedge\ \bar c > c_{\mathrm{full}} \ \wedge\ \tilde r_t \le r_{\mathrm{cap}}. \label{eq:gateB}
\end{align}
Otherwise the meta-review is returned to the author as feedback, the author revises, and the revision is committed as a new state.
Table~\ref{tab:hyper} lists the gate hyperparameters. They are defined per answer format rather than per dataset and were fixed before the experiments reported in this paper.

\begin{table}[!t]
\caption{Acceptance-gate hyperparameters (centralized mode)}
\label{tab:hyper}
\centering
\footnotesize
\setlength{\tabcolsep}{3pt}
\begin{tabular}{lccccc}
\toprule
Task family & $\theta_0$ & $\kappa_{\min}$ & $c_{\min}$ & $c_{\mathrm{full}}$ & $r_{\mathrm{cap}}$ \\
\midrule
\makecell[l]{GSM8K, SVAMP,\\GSM-Hard, StrategyQA} & 0.28 & 2 & 0.65 & 0.80 & 0.45 \\
CIAR & 0.35 & 1 & 0.60 & 0.75 & 0.50 \\
MMLU subsets & 0.38 & 1 & 0.58 & 0.72 & 0.52 \\
\bottomrule
\end{tabular}
\end{table}

\subsection{Overall Procedure}

Algorithm~\ref{alg:wm-agent} summarizes one run. The same procedure covers both architectures; they differ in where the state lives and in the acceptance test of line~\ref{line:accept}.

\begin{algorithm}[t]
\caption{World-Model--Guided Multi-Agent Reasoning}
\label{alg:wm-agent}
\begin{algorithmic}[1]
\REQUIRE Query \(x\), reviewers \(k\), maximum review rounds \(N\)
\ENSURE Final answer \(y\)
\STATE \(s \leftarrow T(\varnothing, \langle\textsc{init}, x\rangle)\)
\STATE Author drafts candidate \(c\); \(a \leftarrow \langle\textsc{draft}, c\rangle\); \(s \leftarrow T(s, a)\)
\IF{\(V(s,a)\) fails}
    \STATE Author re-submits \(c\) given the validator hint; \(s \leftarrow T(s, \langle\textsc{revalidated}, c\rangle)\)
\ENDIF
\FOR{\(n = 1,\ldots,N\)}
    \STATE Each reviewer evaluates \(c\) given \(\sigma(s)\) \(\rightarrow\) reviews \(\mathcal{F}\)
    \STATE \(r \leftarrow R(s, \mathcal{F})\) \COMMENT{trajectory risk, Eq.~\eqref{eq:risk}}
    \STATE Meta-reviewer decides \(d\in\{\text{right},\text{wrong}\}\) given \(c, \mathcal{F}, r\)
    \IF{\(d=\text{right}\) \AND \textsc{Accept}\((r,\mathcal{F})\)} \label{line:accept}
        \RETURN answer of \(c\)
    \ENDIF
    \IF{\(n<N\)}
        \STATE Author revises \(c\) given the meta-review; \(a \leftarrow \langle\textsc{revision}, c\rangle\); \(s \leftarrow T(s,a,n)\)
        \STATE Re-submit once if \(V(s,a)\) fails
    \ENDIF
\ENDFOR
\RETURN answer of the last candidate \(c\)
\end{algorithmic}
\vspace{2pt}
{\footnotesize \textsc{Accept} is Eqs.~\eqref{eq:gateA}--\eqref{eq:gateB} in centralized mode and always true in peer mode.}
\end{algorithm}

\subsection{Coordination Uncertainty Index}
\label{sec:cui}

To characterize when each architecture is preferable we define the \emph{Coordination Uncertainty Index} (CUI), a dataset-level statistic computed from the logs of the first review round.
It is inspired by information-theoretic notions of uncertainty~\cite{shannon1948mathematical} and aggregates four signals of how hard it is for the agents to agree on a stable shared state.
For a dataset \(D\) with \(N\) tasks:

\begin{itemize}
\item \textbf{Reviewer belief entropy.} Let \(p^{(i)}\) be the fraction of reviewers whose decision on task \(i\) is ``right''. Then
\(H_{\mathrm{rev}} = \tfrac1N\sum_i \big[-p^{(i)}\log p^{(i)} - (1-p^{(i)})\log(1-p^{(i)})\big]\).
\item \textbf{Confidence dispersion.} \(V_{\mathrm{conf}} = \tfrac1N\sum_i \mathrm{Var}(c^{(i)}_1,\ldots,c^{(i)}_k)\), the variance of the reviewers' reliability scores.
\item \textbf{Risk instability.} \(V_{\mathrm{risk}} = \tfrac1N\sum_i \mathrm{Var}(r^{(i)}_1,\ldots,r^{(i)}_{n_i})\), the variance of the trajectory risk (Eq.~\eqref{eq:risk}) over the \(n_i\) review rounds of task \(i\) (zero for single-round tasks).
\item \textbf{Answer entropy.} \(H_{\mathrm{ans}}\) is the mean Shannon entropy of the multiset of normalized final answers proposed for a task by the author, the reviewers, and the meta-reviewer.
\end{itemize}

Each component is min--max normalized over the pool of datasets under study (denoted by a hat), and the index is their unweighted mean:
\begin{equation}
\label{eq:cui}
U(D) = \tfrac14\big(\hat H_{\mathrm{rev}} + \hat V_{\mathrm{conf}} + \hat V_{\mathrm{risk}} + \hat H_{\mathrm{ans}}\big).
\end{equation}
We use equal weights to avoid fitting additional free parameters on a small number of datasets.
Because of the normalization, \(U(D)\) is a \emph{relative} index: it ranks the datasets of a pool by coordination uncertainty, and a threshold \(\tau\) is meaningful only with respect to that pool.
Unless stated otherwise, the CUI is computed from peer-mode logs, which are the cheapest to obtain.
